\scanpage{22}
% source lines 1--4
3) $G$ schéma en groupes réductif sur $S$, $B$ sous-schéma en groupes \unsure{simple|lisse}$/S$ induisant sur chaque fibre un Borel. Est-ce que $B$ est un Borel ?? / Oui, car on se ramène au cas semi-simple adjoint\unsure{....}

% source lines 5--9; outlined argument cancelled diagonally.
\deleted{Pour démontrer, on peut supposer que $B$ contient un tore maximal de $G$ \inserted{déployé}. Alors $\unsure{\mathscr L|\mathcal L}=\unsure{\mathfrak t|t}+\sum_{\alpha\in\unsure{P|\Phi}}\unsure{\mathfrak g|\mathscr G}_\alpha$, \unsure{le} $P$ un ``ens de racines positives''. Donc il existe un Borel $B'$ ayant même $\mathscr L$.}

% source lines 10--14
Laissons tomber l'hypothèse $G$ réductif,

\deleted{On peut \unsure{espérer} \unsure{par} \unsure{dessus} que} on se demande si deux tels $B,B'$ sont \inserted{loc\textsuperscript{t}} \uline{conjugués} \deleted{par extension \unsure{fid\textsuperscript{t}}} \inserted{au sens de la top.} fid\textsuperscript{t} plate quasi cpt. Cela se vérifie sur les fibres géométriques : la question
% source line 15
\[
H^{\unsure{1}}(\unsure{B|\mathscr L},\mathfrak g/\mathscr L)=0\qquad??
\]
% source line 16
du point \uline{de vue infinitésimal}.

% source lines 17--19
On peut supposer que $B$ et $B'$ contiennent un $=$ tore $T$, et [\deleted{\unsure{de}} \unsure{quitte} au voisinage du point $s\in S$] que $B_s=B'_s$
% archive sheet number: 21
